\section{Data: Coverage, Noise and the Accuracy Any Surrogate Can Reach}
\label{sec:eda-ns}
\label{sec:eda-bh}

After filtering (Section~\ref{sec:neutron-stars}), the neutron-star table holds $\nsEdaRows$ configurations on $\nsEdaCurves$ curves, one per $(\beta, \lambda)$ pair; the black-hole table holds $\bhEdaRows$ configurations on $\bhEdaCurves$ curves, one per $\beta$, each sampled densely in the horizon radius $\rh$. The analysis leads to six decisions, argued below for both maps with their evidence: the input scale, the target, the role of the parameters, the domain, the split, and the bound the spacing between curves sets on the reachable error. The black-hole answers differ from the neutron stars' in three places (2, 3 and 4), and its data allow far more precision (6).

\paragraph{1.\ The coordinate enters on its natural scale, standardized.}
For neutron stars the central density enters as $\log_{10}\rhoc$: it is positive and multiplicative, and the mass follows $\log_{10}\rhoc$ as one smooth S-shaped curve per $(\beta, \lambda)$ pair\pickref{fig:ns-mass-density}{ (}{)}. All curves share one grid of $\nsEdaDensityGridValues$ densities over $\nsEdaRhocOrders$ orders of magnitude, spaced about evenly in $\rhoc$, so the marginal is a plateau that falls off as curves end at their maximum mass. The representation screen confirms the choice at equal model: on its raw scale $\rhoc$ costs every model between $0.14$ and $0.20$ significant figures on the mass (Section~\ref{sec:representation}).
For black holes $\rh$ spans only $\bhEdaRhOrders$ orders of magnitude, from $\bhEdaRhMin$ to $\bhEdaRhMax$, on a uniform grid of step $\bhEdaRhStep$, with a flat distribution of skew $\bhEdaRhRawSkew$; a logarithm would only make it skewed, so $\rh$ enters as it is. Either input is standardized on the training folds only.

\paragraph{2.\ The charge target is $\log_{10}(\Dch/M)$ for neutron stars, not for black holes.}
The neutron-star charge spans $\nsEdaDOrders$ orders of magnitude, so its target is logarithmic; dividing by the mass makes it a better one (\cref{tab:ns-charge-correlation}). Two coefficients measure how a target depends on another quantity. Pearson's $r$ measures how closely the points follow a straight line, and $r^2$ is the share of the target's variance that a straight line in that quantity explains. Spearman's $\rho$ is Pearson's $r$ computed on ranks: it asks only whether a larger value of one goes with a larger value of the other, whatever the shape of the relation.
Against the mass, $\log_{10}\Dch$ has $\rho = \nsEdaCorrMLogD$: heavier stars tend to carry more charge, so part of the charge merely repeats the mass. For $\log_{10}(\Dch/M)$, $\rho$ falls to $\nsEdaCorrMLogDM$ (and $r$ from $\nsEdaPearsonMLogD$ to $\nsEdaPearsonMLogDM$), so the charge output no longer duplicates the mass output and the network learns the two separately.
Against $\lambda$, $r$ rises from $\nsEdaCorrLambdaLogD$ to $\nsEdaCorrLambdaLogDM$: a straight line in $\lambda$ explains $\nsEdaLambdaLinearPercentLogDM\%$ of the variance of $\log_{10}(\Dch/M)$ against $\nsEdaLambdaLinearPercentLogD\%$ of $\log_{10}\Dch$. Against $\beta$, both targets stay near $r = \nsEdaCorrBetaLogDM$: the coupling has no overall straight-line effect, yet at fixed $\lambda$ it explains much of the variance (decision~3), an interaction the network has to capture.
\review{physical reading of $\Dch/M$ as the dimensionless charge-to-mass ratio.}
For black holes the mass ranges from $\bhEdaMMin$ to $\bhEdaMMax\,\Msun$ and follows the horizon radius almost exactly, with Pearson $r = \bhEdaPearsonMRh$, one nearly straight line per $\beta$ curve (\cref{fig:bh-mass-radius}); it is bounded and of order one, so the mass target is $M$ itself, as for neutron stars. The dilaton shows only near the existence edge: there the curves leave the general-relativistic relation $M = \rh/2$ and fan out by $\beta$, while for larger $\rh$ they merge into one line. Dividing the charge by the mass brings no improvement, because the charge is already almost a function of the mass alone: the Spearman correlation with $M$ is $\bhEdaSpearmanMLogD$ for $\log_{10}\Dch$ and $\bhEdaSpearmanMLogDM$ for $\log_{10}(\Dch/M)$. The target is fixed with the decisions in Section~\ref{sec:preprocessing}.

\InputIfFileExists{91_neutron_stars_tab_charge_correlation}{}{}
\FloatBarrier

\paragraph{3.\ The parameters stay inputs.}
For neutron stars $\lambda$ carries the strongest linear signal in the charge, and at fixed $\lambda$ the coupling $\beta$ still explains $\nsEdaBetaSharePercent\%$ of the variance of $\log_{10}(\Dch/M)$ (median over $\lambda$).
\review{$\beta$ as the first-order coupling and $\lambda$ as a second-order self-interaction, against the field equations (Section~\ref{sec:action}).}
For black holes, at fixed $\rh$ changing $\beta$ moves the mass by only $\bhEdaBetaEffectPercentM\%$ of its spread and $\log_{10}\Dch$ by $\bhEdaBetaEffectPercentLD\%$; $\beta$ sets where each curve starts, not what it holds. Its role is the existence edge: the smallest horizon radius rises from $\bhEdaEdgeRhMinLow$ at $\beta = \bhEdaEdgeBetaLow$ to $\bhEdaEdgeRhMinHigh$ at $\beta = \bhEdaEdgeBetaHigh$ (\cref{fig:validity-domain}). $\beta$ therefore stays an input.
\review{whether the black-hole and neutron-star datasets come from the same theory; $\beta$ spans $\bhEdaEdgeBetaLow$ to $\bhEdaEdgeBetaHigh$ here and $0.4$ to $49.4$ for neutron stars, and matters far less here.}

\paragraph{4.\ The domain is triangular, and the neutron-star data pile up near the maximum mass.}
Only $\nsEdaGridFillPercent\%$ of the $\nsEdaGridCells$ $(\beta, \lambda)$ cells hold a converged curve: the corner of large $\beta$ and large $\lambda$ is empty, and the maximum mass falls towards it, from $\nsEdaMassMax\,\Msun$ at small $\lambda$. Curves also thin out towards that edge, from $\nsEdaRowsPerCurveMax$ rows to $\nsEdaRowsPerCurveMin$. Along each curve the densities are spaced about evenly while $dM/d\rhoc \to 0$ at the turning point, so $\nsEdaPileUpPercent\%$ of all rows lie within $\nsEdaPileUpWindow\,\Msun$ of their curve's maximum: a sampling artefact that weights the region of Hypothesis~$H_2$ heavily. The black-hole domain of validity is the region above the existence edge of decision~3.
\review{$\nsEdaMassMax\,\Msun$ against the heaviest observed neutron stars (about $2.08\,\Msun$, PSR J0740+6620; about $2.35\,\Msun$, PSR J0952$-$0607) and theoretical maxima of about $2.2$ to $2.5\,\Msun$; the single polytropic equation of state sets the dataset's maximum; masses assumed in solar units; add citations.}

\paragraph{5.\ The split groups whole curves.}
The rows of a curve are each other's nearest neighbours: for $\nsEdaSameCurvePercent\%$ of the neutron-star rows the nearest row in standardized $(\beta, \lambda, \log_{10}\rhoc)$ lies on the same curve, at a median distance of $\nsEdaWithinMedian$ against $\nsEdaAcrossMedian$ to another curve. A random split of rows therefore leaks; a one-nearest-neighbour probe on $\log_{10}(\Dch/M)$ measures how much (\cref{tab:ns-split-strategies}, $\nsEdaFolds$ folds). Holding out rows at random gives an error of $\nsEdaSplitRandomRowsDM$, holding out whole curves $\nsEdaSplitCurvesDM$, the honest test of prediction at unseen $(\beta, \lambda)$ pairs. Whole $\beta$ lines cost the same ($\nsEdaSplitBetaLinesDM$), whole $\lambda$ lines are harder ($\nsEdaSplitLambdaLinesDM$), and the rim predicted from the interior is hardest ($\nsEdaSplitRimDM$), so extrapolation is reported separately and not claimed. For the mass the order reverses ($\nsEdaSplitRandomRowsM$ against $\nsEdaSplitCurvesM$): the mass depends mostly on $\rhoc$, and a neighbouring curve at the same density predicts it better than the next density on the same curve.
Each black-hole curve is so dense that every row's nearest neighbour lies on its own curve ($\bhEdaSameCurvePercent\%$; median distance $\bhEdaWithinMedian$ against $\bhEdaAcrossMedian$ to the next curve), so there too a random split leaks and the split groups whole $\beta$ curves. The probe (\cref{tab:bh-split-strategies}, $\bhEdaFolds$ folds) shows it for the charge, which $\beta$ affects most: $\log_{10}\Dch$ scores $\bhEdaSplitRandomRowsLD$ with rows held out at random against $\bhEdaSplitCurvesLD$ with whole curves. Holding out the outermost curves scores as holding out any curve ($\bhEdaSplitOuterCurvesLD$), so extrapolating in $\beta$ is no harder than interpolating; the mass alone reverses the order ($\bhEdaSplitRandomRowsM$ against $\bhEdaSplitCurvesM$), because all curves share one $\rh$ grid.

\InputIfFileExists{91_neutron_stars_tab_split_strategies}{}{}
\FloatBarrier

\paragraph{6.\ The data bound the reachable error.}
How closely can any surrogate reproduce a curve it has not seen? Predicting each interior curve from the two curves on each side -- a cubic spline in $\beta$ at fixed $\lambda$, or in $\lambda$ at fixed $\beta$, after interpolating each neighbour along $\log_{10}\rhoc$ -- gives the error a smooth interpolation of the grid already reaches, the scale against which the models of Section~\ref{sec:results} are judged (\cref{tab:ns-ceilings}, \cref{fig:ns-ceilings}). The data themselves are not the limit: the noise along a curve is $\nsEdaCeilMNoise$ for the mass and $\nsEdaCeilDNoise$ for the charge (median), and a spline through every other row of a curve predicts the rows between to $\nsEdaCeilMAlong$ and $\nsEdaCeilDAlong$ (95th percentile). The limit is the spacing between curves. Across them the mass keeps its error below $\nsEdaCeilMAcrossLambda$ in $\lambda$ and $\nsEdaCeilMAcrossBeta$ in $\beta$ at the 95th percentile, about $\nsEdaCeilMFigures$ significant figures; the charge reaches only $\nsEdaCeilDAcrossLambda$ and $\nsEdaCeilDAcrossBeta$, about $\nsEdaCeilDFigures$ figures. Across $\beta$ the charge's error is heavy-tailed: its median stays near five figures while its 95th percentile does not.\review{Physical reading: are the weakest charges the onset of scalarization, where $\Dch$ grows fastest with the couplings?}
As an uncertainty of the predicted values this bound is a floor, about $\nsEdaCeilMSigma\,\Msun$ on the mass and $\nsEdaCeilDSigma$ on the charge (95th percentiles, as the errors are heavy-tailed), under the error of any surrogate trained on this grid. A neutron-star surrogate is therefore judged against about four significant figures for the mass and two to three for the charge: a bound the data set, not a property of any model.
The same test on the black holes -- each interior curve predicted from the two curves on each side, a cubic spline in $\beta$ after interpolating each neighbour along $\rh$ -- gives the error the $\beta$ grid reaches (\cref{tab:bh-ceilings}, \cref{fig:bh-ceilings}), far less than for neutron stars. The noise along a curve is $\bhEdaCeilMNoise$ for the mass and $\bhEdaCeilDNoise$ for the charge (median), a spline through every other row predicts the rows between to $\bhEdaCeilMAlong$ and $\bhEdaCeilDAlong$, and across curves the 95th percentile of the error is $\bhEdaCeilMAcrossBeta$ for the mass and $\bhEdaCeilDAcrossBeta$ for the charge: about $\bhEdaCeilMFigures$ and $\bhEdaCeilDFigures$ significant figures. As an uncertainty of the predicted values, at the median mass $M = \bhEdaCeilMAt\,\Msun$ and the median charge $\Dch = \bhEdaCeilDAt$ this bound is $\bhEdaCeilMSigma\,\Msun$ and $\bhEdaCeilDSigma$, or $\bhEdaCeilDSigmaDex$ in $\log_{10}\Dch$, at the 95th percentile: a floor on a surrogate's error, to be set beside the observational uncertainty in a future MCMC inference. The error is not spread evenly: it rises by orders of magnitude at the smallest radii, where each curve starts at its own existence edge (decision~3), and stays near the noise everywhere else. Even with only $\bhEdaCurves$ curves, interpolation between them reaches this bound; whether a surrogate keeps it between curves it was not trained on, and near the existence edge, is the black-hole question for Section~\ref{sec:results}.

\InputIfFileExists{91_neutron_stars_tab_ceilings}{}{}
\FloatBarrier
\InputIfFileExists{92_black_holes_tab_split_strategies}{}{}
\FloatBarrier
\InputIfFileExists{92_black_holes_tab_ceilings}{}{}
\FloatBarrier
